%%%%%%%%%%%%%%%%%%%%%%%%%%%%%%%%%%%%%%%%%%%%%%%%%%%%%%%%%
%%%%%%%%%%%%%%%%         EXPERIMENTACION              %%%%%%%%%%%%%%%%%%%%%
%%%%%%%%%%%%%%%%%%%%%%%%%%%%%%%%%%%%%%%%%%%%%%%%%%%%%%%%%


\chapter{Conclusiones}
\label{chap:conclussions}

\chapterintro{
En este capítulo se presentan las conclusiones extraídas tras el desarrollo de la presente Tesis Doctoral. Las conclusiones se basan en los tres artículos presentados, así como en la investigación transversal realizada
}

Los resultados obtenidos de esta Tesis Doctoral han permitido formular las siguientes conclusiones:

\begin{enumerate}

\item En el trabajo \cite{paper1}, a partir de una formada por 79 pacientes con riesgo de sufrir DT2, se demuestra que, a nivel de filo, es posible identificar señales o indicadores de una patología potencial. En este caso, el ratio de Bacteroidetes/Firmicutes varía a lo largo de la progresión de SM y de DT2. A niveles más específicos, se revela que taxones del género \textit{Blautia}, \textit{Tyzzerella} y \textit{Dorea} están enriquecidos en el perfil microbiómico de pacientes que sufren SM. Por otro lado, se presenta que el perfil microbiómico de pacientes con DT2 se caracteriza por una mayor abundancia de taxones como \textit{Dorea}, \textit{Prevotella} y \textit{Dialister}.

\item En el trabajo \cite{paper2}, se ha obtenido una firma metagenómica compuesta por 25 especies capaz de identificar a pacientes diagnosticados con DT1. En esta firma destacan, por su importancia, taxones del género \textit{Bacteroides} y \textit{Prevotella}, que pueden servir como biomarcadores en el caso de la patología de DT1. Además, se pone de manifiesto que un modelo de RF entrenado con dichas especies tiene un poder predictivo capaz de clasificar a los pacientes en sanos, seroconvertidos (estadío intermedio) y afectados por DT1.

\item En el trabajo \cite{paper3}, se ha propuesto un nuevo enfoque de estratificación de pacientes con VB a partir de datos metagenómicos. Este subtipado presenta información complementaria cuando se compara con otros métodos de subtipado del estado del arte. Además, en el contexto de VB, cuando se utiliza la información extraída de este método de subtipado, se mejora la predicción de la respuesta al tratamiento con metronidazol.

\item Finalmente, la investigación desarrollada en esta Tesis Doctoral confirma el potencial del microbioma como fuente de biomarcadores clínicos en varias patologías. De manera paralela, se demuestra el potencial de utilizar técnicas de ML para la construcción de modelos predictivos, y la capacidad de dichas técnicas para lidiar con los obstáculos inherentes al tipo de dato que es el microbioma.

\end{enumerate}